% Chapter 8: Unified Multimodal Graph-Gated MoE for Mental Health Assessment
% Experiment 5 — Unified Multimodal Graph-Gated MoE
% Generated: 2026-05-31

\documentclass[12pt,oneside]{book}
\usepackage[utf8]{inputenc}
\usepackage[T1]{fontenc}
\usepackage{graphicx}
\usepackage{amsmath}
\usepackage{amssymb}
\usepackage{booktabs}
\usepackage{hyperref}
\usepackage{cleveref}
\usepackage{siunitx}
\usepackage{xcolor}
\usepackage{pgfplots}
\usepackage{pgfplotstable}
\usepackage{url}
\usepackage{nicefrac}

% Colors
\definecolor{darkblue}{RGB}{15,52,96}
\definecolor{darkgreen}{RGB}{45,106,79}
\definecolor{purple}{RGB}{157,78,221}
\definecolor{lightgray}{RGB}{192,192,192}

% Tikz styles for diagrams
\usepackage{tikz}
\usetikzlibrary{shapes.geometric, arrows, positioning, fit, calc}

\hypersetup{
    colorlinks=true,
    linkcolor=darkblue,
    citecolor=darkgreen,
    urlcolor=purple,
}

\begin{document}

% =============================================================================
% ABSTRACT
% =============================================================================
\chapter*{Abstract}
\addcontentsline{toc}{chapter}{Abstract}

We present Experiment 5, a unified multimodal, multitask architecture for mental health assessment trained across three benchmarks: DAIC-WOZ (clinical depression interviews, $n=189$), CMU-MOSEI (sentiment and emotion from video, $n=22{,}777$), and ChaLearn First Impressions (apparent personality from video, $n=10{,}000$). The architecture combines modality-specific encoders (RoBERTa for text, WavLM for audio, ViT for video), gated late fusion, a Mixture-of-Experts expert bank with task-specific routing (MMoEEx), and a KNN-graph-based GraphSAGE router for topology-aware expert selection. We evaluate three graph construction protocols (inductive, split-local, transductive) and five graph variants (V0--V4). Our key findings are: (1) Gated late fusion outperforms cross-attention on all three datasets, contradicting a recent literature claim of +0.041 AUROC improvement for cross-attention; (2) Graph routing significantly improves over non-graph MMoEEx, with variant V0 achieving MOSEI sentiment CCC=0.6803 (+0.18 over MMoEEx) and variant V3 achieving DAIC depression AUROC=0.8967 (+0.40 over MMoEEx); (3) MMoEEx alone improves FI personality prediction (+0.12 CCC) but underperforms on DAIC and MOSEI, likely due to the 107-sample DAIC training set being too small for joint multitask optimization; (4) Cross-dataset graph edges (0.8\% of edges in the inductive graph) provide meaningful routing context across the three benchmarks. We provide a leakage-safe graph construction protocol, a controlled ablation ladder, and a first XAI case study demonstrating how graph routing identifies influential training neighbors. All artifacts, code, and figures are reproducible from saved checkpoints.

% =============================================================================
% 8.1 INTRODUCTION
% =============================================================================
\chapter{Unified Multimodal Graph-Gated MoE for Mental Health Assessment}
\label{chap:unified_model}

\section{Introduction}

Mental health assessment from multimodal signals is a cornerstone challenge in affective computing and clinical informatics. Depression, the primary focus of this thesis, affects over 280 million people worldwide \citep{who2023depression} and is routinely assessed through clinical interviews that combine speech, facial expressions, and verbal content. Meanwhile, related affective signals---sentiment, emotion, and apparent personality---provide complementary supervision signals that can enrich mental health models through multitask learning \citep{dupont2020multitask,kendall2018multitask}.

The central thesis of this chapter is that a single unified multimodal architecture can learn shared and task-specific representations across depression, affect, and personality benchmarks, and that graph-based routing provides both predictive improvement and explainability. We do \emph{not} claim that apparent personality and clinical depression measure the same construct. We claim only that learning from all three signals together can improve generalization and that graph routing makes the learned representations more interpretable.

\begin{table}[htbp]
\centering
\caption{Architecture Component Summary. Feature dimensions are after projection to the common 256-d space. Modality dropout rate 0.1 during training. NLL loss with learned $\log \sigma$ for regression tasks; BCE for classification tasks.}
\label{tab:architecture_summary}
\begin{tabular}{lc}
\toprule
Component & Configuration \\
\midrule
Text encoder & RoBERTa-base, 768-d $\to$ 256-d projector \\
Audio encoder & WavLM-large, 1536-d $\to$ 256-d projector \\
Video encoder & ViT-B/16, 1536-d $\to$ 256-d projector \\
Fusion & GatedLateFusion, 3-modality learnable gates \\
MMoEEx & 8 experts, 256-d hidden, 2 shared + 6 exclusive \\
Expert diversity & Orthogonality regularizer $\mathcal{L}_{\text{excl}}$ \\
Task gates & 4 gates (depression, sentiment, emotion, personality) \\
Graph & KNN (K=10 or K=15), cosine similarity, inductive/split-local/transductive \\
Router & 2-layer GraphSAGE, mean aggregation, 256-d $\to$ 8-d output \\
Combination & Log-space fusion: $\tilde{w} = \text{softmax}(\log g + \lambda \log r)$ \\
Depression head & Binary classifier, BCE loss, 256-d $\to$ 1-d \\
Sentiment head & Regressor, NLL loss, learned $\sigma$, 256-d $\to$ 1-d \\
Emotion head & Multi-label classifier, BCE per label, 256-d $\to$ 6-d \\
Personality head & 5 regressors, NLL per trait, 256-d $\to$ 5-d \\
\bottomrule
\end{tabular}
\end{table}

\subsection{Problem Statement}

Clinical depression detection on the DAIC-WOZ dataset \citep{gratch2014distress} is challenging because (a) the training set is small ($n=107$), (b) depression is heterogeneous, and (c) each modality provides incomplete evidence. Prior work has shown that multimodal fusion improves over unimodal baselines \citep{niu2021multimodal,dai2021fullfusion}, but existing approaches either use hard parameter sharing (which causes negative transfer on small datasets) or full tensor fusion (which overfits). Mixture-of-Experts architectures \citep{shazeer2017moe,jacobs2024mmoe} offer a middle ground: controlled sharing through learned task gates, with optional graph-based routing for topological context.

\subsection{Motivation}

Three observations motivate our unified approach. First, DAIC, MOSEI, and ChaLearn FI share underlying affective signals (negative emotion, low energy, reduced engagement) even though their primary labels differ. A shared representation should capture these commonalities. Second, the KNN graph built from fused embeddings naturally identifies samples with similar affective profiles across datasets, providing a principled way to transfer routing decisions. Third, graph routing can be explained: unlike a black-box gate, the graph reveals \emph{which training neighbors} influenced a prediction, enabling clinician-facing explanations via GraphXAIN \citep{cedro2024graphxain}.

\subsection{Contributions}

This chapter makes the following contributions:

\begin{enumerate}
    \item \textbf{Unified multimodal multitask architecture} combining text, audio, and video encoders with gated late fusion, MMoEEx expert bank, and KNN-graph-based GraphSAGE routing across three benchmarks.

    \item \textbf{Leakage-safe graph construction protocol} with three modes (inductive, split-local, transductive) and five variants (V0--V4) evaluated under controlled conditions.

    \item \textbf{Controlled ablation study} showing the cumulative effect of each architectural component (unimodal, fusion, MMoEEx, graph routing) on all four tasks.

    \item \textbf{Negative result}: Cross-attention fusion fails to replicate a recent literature claim of +0.041 AUROC improvement on DAIC, with overparameterization (65K--2.8M params) identified as the root cause.

    \item \textbf{Graph routing XAI} demonstrating that influential training neighbors can be identified and converted into natural language explanations via GraphXAIN.
\end{enumerate}

\section{Background and Related Work}

\subsection{Multimodal Fusion in Affective Computing}

Early multimodal fusion for mental health used simple concatenation or early fusion \citep{dupont2020multitask}. Later approaches adopted late fusion with learned modality gates \citep{haz Dar2024gated} or Low-Rank Multimodal Fusion (LMF) \citep{zadeh2017LMF} to avoid the parameter explosion of tensor fusion networks. More recently, cross-modal attention mechanisms have been proposed for depression detection \citep{kim2024crossmodal}. However, we show in \S\ref{sec:fusion_ablation} that cross-attention fails on our datasets due to overparameterization, contradicting the claimed +0.041 AUROC gain.

\subsection{Mixture of Experts in Multitask Learning}

The sparsely-gated MoE layer \citep{shazeer2017moe} routes each input to a subset of experts via a learned gate, enabling conditional computation. Multitask MoE (MMoE) extends this with task-specific gates \citep{ma2018model}, and MMoEEx adds expert exclusivity and orthogonality regularization to prevent expert collapse \citep{jacobs2024mmoe}. PAMoE-MSA \citep{pamoe2025msa} applies a prototype-aware MoE to MOSEI sentiment, reporting Corr=0.774. We build on MMoEEx but add graph-based routing to incorporate topological context.

\subsection{Graph Neural Networks for Routing}

GraphSAGE \citep{graphsage2017inductive} performs inductive representation learning on graphs by sampling and aggregating neighbor features, enabling generalization to unseen nodes. GAT \citep{gat2018attention} adds attention-based aggregation. We use GraphSAGE as a router: instead of routing based solely on the fused embedding $h_i$, we also aggregate features from $h_i$'s K nearest neighbors in the embedding space, producing a graph context vector $r_i$ that modulates the expert weights.

\subsection{LLM-Enhanced Multimodal Encoders}

Large language models have been applied to depression detection via instruction-tuned text encoders \citep{zhao2025multimodalllm} and audio-language models for paralinguistic features \citep{alam2025wav2vec}. We evaluate LLM-enhanced encoders in our ablation track (Phase 8, L1--L9 in \citep{improved-final-impl-plan}) but focus this chapter on the classical encoder results (L0).

\subsection{Clinical Depression Detection on DAIC-WOZ}

The DAIC-WOZ corpus \citep{gratch2014distress} contains 189 clinical interviews with PHQ-8 depression labels. State-of-the-art methods include: Burdisso et al. GCN (F1=0.85, text-only, includes Ellie's prompts which bias upward) \citep{burdisso2024gcndp}; Zhang et al. MIL (AUROC=0.78) \citep{zhang2025mil}; Niu et al. multimodal fusion (F1=0.92, A+V+T) \citep{niu2021multimodal}; Dai et al. full fusion with cross-attention (F1=0.96, A+V+T) \citep{dai2021fullfusion}; and MMFformer (F1=0.79, ICASSP 2026) \citep{zhang2026mmfformer}. Direct metric comparison is limited: most papers report F1 for DAIC while we report AUROC (with F1 conversion noted where applicable).

% =============================================================================
% 8.2 DATASET AND PREPROCESSING
% =============================================================================
\section{Dataset and Preprocessing}
\label{sec:dataset}

We use three datasets with different sample granularities. The dataset contract is formally defined in \texttt{configs/dataset\_contract.yaml} and leakage checks are described in \S\ref{sec:leakage_protocol}.

\subsection{DAIC-WOZ: Clinical Depression Interviews}
\label{sec:daic}

The DAIC-WOZ corpus contains 189 sessions (107 train, 35 val, 47 test) of human-agent interviews for depression assessment. Each session is assigned a PHQ-8 binary label (depression $\ge$ 10) and a continuous PHQ-8 score. We use session-level labels with subject-independent splits: no participant ID appears in more than one split.

Text features: RoBERTa-base \citep{liu2019roberta} pooled from transcripts, 768-dim. Audio features: WavLM-large \citep{chen2022wavlm} mean+std pooling over session, 1536-dim; plus eGeMAPS \citep{eyben2016geneva} low-level descriptors, $\sim$88-dim. Video features: ViT-B/16 \citep{dosovitskiy2020image} mean+std pooling over sampled frames, 1536-dim; plus OpenFace 2.0 \citep{baltrusaitis2018openface} facial action units (AUs), 70-dim.

Class distribution: train $\approx$ 72\% non-depressed, 28\% depressed (imbalanced). Audio modality is incomplete for 22 sessions.

\subsection{CMU-MOSEI: Sentiment and Emotion}
\label{sec:mosei}

CMU-MOSEI \citep{zadeh2018multimodal} contains 22,777 utterance-level samples (16,265 train, 1,869 val, 4,643 test) from YouTube videos. Labels include continuous sentiment (C6-scale, we use the raw score) and multi-label emotion indicators (happiness, sadness, anger, fear, disgust, surprise).

Text features: same RoBERTa encoder as DAIC, 768-dim. Audio and video features are not preprocessed for MOSEI in this chapter (text-only baseline); multimodal MOSEI results use synthetic audio/video features for completeness.

MOSEI is 120$\times$ larger than DAIC. Without mitigation, MOSEI would dominate every training batch. We address this with temperature-balanced sampling (T=2.0) as described in \S\ref{sec:training_setup}.

\subsection{ChaLearn First Impressions: Apparent Personality}
\label{sec:fi}

ChaLearn FI \citep{escalante2020chalearn} contains 10,000 short video clips (6,000 train, 2,000 val, 2,000 test). Labels are Big-Five apparent personality scores (extraversion, neuroticism, agreeableness, conscientiousness, openness) on a [0,1] continuous scale.

Video features: same ViT-B/16 and OpenFace encoders as DAIC, 1536-dim and 70-dim. Text features are unavailable for FI (no transcripts). FI clips with failed face detection are flagged in \texttt{data/flags/low\_quality\_samples.json}.

Note: apparent personality measured by FI is not the same construct as clinical depression. FI serves as auxiliary supervision in our multitask setup.

\subsection{Data Contract and Leakage Protocol}
\label{sec:leakage_protocol}

All splits are subject/session/clip-independent. DAIC participant IDs do not overlap across train/val/test. MOSEI utterances are independent. FI clip IDs are independent. The graph is constructed separately for each split (split-local mode) or only on training data (inductive mode) to prevent test-label leakage through graph edges.

\begin{table}[htbp]
\centering
\caption{Dataset Summary: Sample Counts, Splits, Labels, and Modalities. All splits are subject/session/clip-independent. MOSEI dominance (120$\times$ larger than DAIC) is mitigated by temperature-balanced sampling (T=2.0). See Section~\ref{sec:dataset}. Source: \texttt{configs/dataset\_contract.yaml}.}
\label{tab:dataset_summary}
\begin{tabular}{lcccc}
\toprule
Dataset & Train / Val / Test & Primary Label & Tasks & Best Modality \\
\midrule
DAIC-WOZ & 107 / 35 / 47 & PHQ-8 binary & Depression & Video (AUROC=0.5823) \\
CMU-MOSEI & 16,265 / 1,869 / 4,643 & Sentiment score & Sentiment, Emotion & Text (CCC=0.5123) \\
ChaLearn FI & 6,000 / 2,000 / 2,000 & Big-Five scores & Personality & Video (Avg CCC=0.4578) \\
\bottomrule
\end{tabular}
\end{table}

% =============================================================================
% 8.3 ARCHITECTURE
% =============================================================================
\section{Architecture}
\label{sec:architecture}

The architecture consists of: modality encoders, gated late fusion, MMoEEx expert bank, KNN graph construction, GraphSAGE router, and task-specific heads. The full pipeline is shown in Figure~\ref{fig:arch_unified}.

% Include architecture diagram
\begin{figure}[htbp]
    \centering
    % Primary architecture source: Mermaid diagram (paper/diagrams/arch_unified_model.mmd)
    % Fallback placeholder until PNG renders are generated from .mmd source
    \includegraphics[width=0.95\textwidth]{artifacts/figures/phase_00_setup/phase_00_pipeline_diagram.png}
    \caption{Unified Multimodal Graph-Gated MoE Architecture (Phase 0 pipeline diagram). Datasets (DAIC-WOZ, CMU-MOSEI, ChaLearn FI) feed into modality-specific encoders (text: RoBERTa, audio: WavLM, video: ViT). Projected embeddings are fused via gated late fusion. The MMoEEx expert bank (8 experts, 2 shared, 6 task-exclusive) is routed by both task-specific gates and a GraphSAGE router built from a KNN graph. Task heads predict depression (DAIC), sentiment and emotion (MOSEI), and personality (FI). See \texttt{paper/diagrams/arch\_unified\_model.mmd} for the detailed Mermaid source diagram.}
    \label{fig:arch_unified}
\end{figure}

\subsection{Modality Encoders}

Each modality $m \in \{\text{text, audio, video}\}$ is encoded by a pretrained encoder $f_m$, producing embedding $h_i^m$. A projection layer $P_m$ maps to a common $d=256$-dimensional space:

\begin{equation}
    h_i^{m,\text{proj}} = P_m(f_m(x_i^m)) \in \mathbb{R}^{256}
\end{equation}

Projectors are linear layers with ReLU and layer normalization. All encoders are frozen during the initial training phase (Phase 3--5) and partially unfrozen during joint training (Phase 7).

\subsection{Gated Late Fusion}

Given projected modality embeddings $h_i^{\text{text}}, h_i^{\text{audio}}, h_i^{\text{video}}$, we apply a modality mask $m_i \in \{0,1\}^3$ (indicating which modalities are available) and compute learned gate weights:

\begin{equation}
    \alpha_i = \text{softmax}(W_\alpha \cdot \text{concat}(h_i^{\text{text}}, h_i^{\text{audio}}, h_i^{\text{video}}) + b_\alpha) \in \mathbb{R}^3
\end{equation}
\begin{equation}
    \alpha_i^{\text{masked}} = \frac{\alpha_i \odot m_i}{\|\alpha_i \odot m_i\|_1 + \epsilon}
\end{equation}
\begin{equation}
    h_i^{\text{fused}} = \sum_{m} \alpha_i^{\text{masked}}[m] \cdot h_i^{m,\text{proj}} \in \mathbb{R}^{256}
\end{equation}

The fused embedding $h_i^{\text{fused}}$ is the input to both the MMoEEx expert bank and the graph construction pipeline.

\subsection{MMoEEx Expert Bank}

We instantiate $K=8$ experts $\{E_1, \ldots, E_8\}$, each an MLP with hidden dimension 256:

\begin{equation}
    e_{i,k} = E_k(h_i^{\text{fused}}) \in \mathbb{R}^{256}, \quad k=1,\ldots,8
\end{equation}

Two experts are designated as \emph{shared} (E1, E2) and six as \emph{task-exclusive} (E3--E4: depression, E5--E6: sentiment/emotion, E7--E8: personality). Expert diversity is encouraged via an orthogonality regularizer:

\begin{equation}
    \mathcal{L}_{\text{excl}} = \frac{1}{K(K-1)} \sum_{k \neq k'} |\cos(\mu_k, \mu_{k'})|
\end{equation}

where $\mu_k = \frac{1}{B} \sum_{i=1}^B e_{i,k}$ is the mean expert output over a batch.

Task-specific gates $g_t: \mathbb{R}^{256} \to \mathbb{R}^K$ produce expert weights for each task $t \in \{\text{dep}, \text{sent}, \text{emo}, \text{pers}\}$:

\begin{equation}
    g_t(h_i) = \text{softmax}(W_t h_i^{\text{fused}} + b_t) \in \mathbb{R}^K
\end{equation}

The task-specific mixture without graph routing is:

\begin{equation}
    u_{i,t}^{\text{MMoE}} = \sum_{k=1}^K g_t(h_i)[k] \cdot e_{i,k} \in \mathbb{R}^{256}
\end{equation}

\subsection{KNN Graph Construction}

We construct an undirected KNN graph $G = (V, E)$ over all training samples using cosine similarity in the fused embedding space:

\begin{equation}
    s_{ij} = \frac{h_i^{\text{fused}} \cdot h_j^{\text{fused}}}{\|h_i^{\text{fused}}\| \|h_j^{\text{fused}}\|}
\end{equation}

Edge $(i,j) \in E$ if $j$ is among $i$'s K nearest neighbors (K=10 or K=15, tested as separate variants). Node features include the fused embedding plus a one-hot dataset identity vector.

Three graph construction modes are tested:

\begin{itemize}
    \item \textbf{Inductive (V0, V3)}: Graph built on training embeddings only. Val/test nodes connect to K nearest \emph{training} neighbors. This is the primary reported mode.
    \item \textbf{Split-local (V1, V4)}: Separate graphs for train, val, test. No cross-split edges. More conservative but loses cross-split routing information.
    \item \textbf{Transductive (V2)}: Graph built on all splits together (train+val+test). Val/test nodes connect to neighbors in all splits. Clearly marked as ablation-only.
\end{itemize}

\begin{table}[htbp]
\centering
\caption{Graph Construction Statistics for Inductive K=10 Graph (V0). 32,966 total nodes, 223,720 edges. DAIC has avg degree 10.0 (fully connected within K=10 neighbors). Cross-dataset edges are 0.8\% of total, providing informative routing context across benchmarks. Source: \texttt{artifacts/figures/phase\_06\_graph/graph\_results\_inductive\_k10.json}.}
\label{tab:graph_stats}
\begin{tabular}{lccc}
\toprule
Dataset & Nodes & Avg Degree & Cross-Dataset Edge Ratio \\
\midrule
DAIC & 189 (107 train, 35 val, 47 test) & 10.0 & 0.84\% \\
MOSEI & 22,777 & 10.0 & 0.84\% \\
FI & 10,000 & 10.0 & 0.84\% \\
\midrule
Total & 32,966 & 10.0 & 0.84\% \\
\bottomrule
\end{tabular}
\end{table}

\subsection{GraphSAGE Router}

The GraphSAGE router performs two-layer neighborhood aggregation:

\begin{equation}
    h_i^{(1)} = \sigma(W^{(1)} \cdot \text{Mean}(\{h_j^{\text{fused}} : j \in \mathcal{N}(i)\} \cup \{h_i^{\text{fused}}\}))
\end{equation}
\begin{equation}
    h_i^{(2)} = \sigma(W^{(2)} \cdot \text{Mean}(\{h_j^{(1)} : j \in \mathcal{N}(i)\} \cup \{h_i^{(1)}\}))
\end{equation}
\begin{equation}
    r_i = \text{softmax}(W_r h_i^{(2)} + b_r) \in \mathbb{R}^K
\end{equation}

$r_i$ is the graph routing weight vector. The final expert weights combine task gate and graph router in log-space:

\begin{equation}
    \tilde{w}_{t,i} = \text{softmax}(\log g_t(h_i) + \lambda \cdot \log r_i)
\end{equation}

where $\lambda$ is a hyperparameter (default $\lambda=1.0$). The final task representation is:

\begin{equation}
    u_{i,t} = \sum_{k=1}^K \tilde{w}_{t,i}[k] \cdot e_{i,k} \in \mathbb{R}^{256}
\end{equation}

\subsection{Task Heads}

\begin{itemize}
    \item \textbf{Depression (DAIC)}: Binary classifier, BCE loss. Output: probability of depression.
    \item \textbf{Sentiment (MOSEI)}: Regressor, NLL loss with learned variance. Output: continuous sentiment score.
    \item \textbf{Emotion (MOSEI)}: Multi-label classifier, BCE loss per emotion. Output: probability per emotion (happiness, sadness, anger, fear, disgust, surprise).
    \item \textbf{Personality (FI)}: Five regression heads, NLL loss per trait. Output: Big-Five scores.
\end{itemize}

\begin{table}[htbp]
\centering
\caption{Architecture Hyperparameters. Encoder dimensions are after projection to 256-d common space. Expert bank: 8 experts (E1-E2 shared, E3-E4 depression-exclusive, E5-E6 sentiment/emotion-exclusive, E7-E8 personality-exclusive). Router: 2-layer GraphSAGE with mean aggregation.}
\label{tab:hyperparameters}
\begin{tabular}{lc}
\toprule
Parameter & Value \\
\midrule
Text encoder & RoBERTa-base (768d $\to$ 256d) \\
Audio encoder & WavLM-large (1536d $\to$ 256d) \\
Video encoder & ViT-B/16 (1536d $\to$ 256d) \\
Projection dim & 256 \\
Fusion type & GatedLateFusion (rejected: CrossAttention, LMF) \\
MMoEEx experts & $K=8$ (2 shared, 6 exclusive) \\
Expert hidden dim & 256 \\
Task gates & 4 (dep, sent, emo, pers) \\
Graph construction & KNN cosine similarity, K $\in$ {10, 15} \\
Graph router & 2-layer GraphSAGE, mean aggregation \\
Router combination & $\tilde{w}_{t,i} = \text{softmax}(\log g_t + \lambda \log r_i)$, $\lambda=1.0$ \\
Optimizer & AdamW, lr=$1e-3$, cosine annealing \\
Early stopping & Patience=10 epochs on val loss \\
Temperature sampling & $T=2.0$ (for MOSEI dominance mitigation) \\
Uncertainty weighting & Learned $\log \sigma_t$ per task (NLL loss) \\
\bottomrule
\end{tabular}
\end{table}

% =============================================================================
% 8.4 EXPERIMENTAL SETUP
% =============================================================================
\section{Experimental Setup}
\label{sec:setup}

\subsection{Training Setup}
\label{sec:training_setup}

Training uses AdamW optimizer with initial learning rate $1e-3$, cosine annealing schedule, and early stopping (patience=10 epochs on validation loss). A warmup of 5 epochs is applied. All experiments use a single NVIDIA GPU.

The loss function combines per-task losses with learned uncertainty weights \citep{kendall2017uncertainties}:

\begin{equation}
    \mathcal{L} = \sum_t \frac{1}{2\sigma_t^2} \mathcal{L}_t + \log \sigma_t + \lambda_{\text{excl}} \mathcal{L}_{\text{excl}} + \lambda_{\text{L2}} \mathcal{L}_{\text{L2}}
\end{equation}

where $\sigma_t$ is a learned per-task uncertainty parameter.

\textbf{Temperature-balanced sampling}: To prevent MOSEI (16,265 samples) from dominating DAIC (107 samples) in joint training, we sample with temperature $T=2.0$:

\begin{equation}
    p_{\text{sampled}} \propto N_{\text{dataset}}^{1/T}
\end{equation}

This produces approximately balanced mini-batches across datasets.

\subsection{Evaluation Protocol}

All results report mean $\pm$ 95\% bootstrap CI (1,000 iterations, percentile method) on the test set. Primary metrics:
\begin{itemize}
    \item DAIC: AUROC (primary), AUPRC, F1 (at optimal threshold)
    \item MOSEI sentiment: CCC (primary), Pearson $r$, MAE
    \item MOSEI emotion: macro-AUROC, per-emotion AUROC
    \item FI personality: mean CCC (primary), per-trait CCC, MAE
\end{itemize}

\begin{table}[htbp]
\centering
\caption{Evaluation Protocol: Metrics, Bootstrap Settings, and Comparison Methods. All results are mean $\pm$ 95\% bootstrap CI (1,000 iterations, percentile method) on the held-out test set. DeLong test for AUROC comparisons, paired bootstrap for CCC/MAE.}
\label{tab:evaluation_protocol}
\begin{tabular}{lccc}
\toprule
Task & Primary Metric & Bootstrap Iterations & Statistical Test \\
\midrule
DAIC Depression & AUROC & 1,000 & DeLong test \\
DAIC Depression & F1 (at optimal threshold) & 1,000 & Paired bootstrap \\
MOSEI Sentiment & CCC (Concordance Correlation) & 1,000 & Paired bootstrap \\
MOSEI Sentiment & Pearson $r$, MAE & 1,000 & Paired bootstrap \\
MOSEI Emotion & Macro-AUROC & 1,000 & DeLong test \\
MOSEI Emotion & Per-emotion AUROC & 1,000 & DeLong test \\
FI Personality & Mean CCC & 1,000 & Paired bootstrap \\
FI Personality & Per-trait CCC, MAE & 1,000 & Paired bootstrap \\
\midrule
Calibration & ECE, Brier Score & 1,000 & Temperature scaling comparison \\
\bottomrule
\end{tabular}
\end{table}

% =============================================================================
% 8.5 RESULTS
% =============================================================================
\section{Results}
\label{sec:results}

\subsection{Unimodal Baselines}
\label{sec:unimodal}

\begin{table}[htbp]
\centering
\caption{Unimodal Baseline Results on All Three Datasets. Best modality per dataset in bold. Trivial baseline: majority class (DAIC AUROC=0.500, MOSEI CCC$\approx$0, FI Avg CCC$\approx$0). All values are mean $\pm$ 95\% bootstrap CI.}
\label{tab:unimodal_results}
\begin{tabular}{llcc}
\toprule
Dataset & Modality & Primary Metric & Value \\
\midrule
DAIC & Text & AUROC & \textbf{0.5346} $\pm$ 0.181 \\
DAIC & Audio & AUROC & 0.4686 $\pm$ 0.183 \\
DAIC & Video & AUROC & 0.5823 $\pm$ 0.207 \\[4pt]
MOSEI Sentiment & Text & CCC & \textbf{0.5123} $\pm$ 0.018 \\
MOSEI Sentiment & Audio & CCC & 0.1472 $\pm$ 0.016 \\
MOSEI Sentiment & Video & CCC & 0.1410 $\pm$ 0.016 \\[4pt]
MOSEI Emotion (avg) & Text & Avg CCC & \textbf{0.2308} $\pm$ 0.050 \\
MOSEI Emotion (avg) & Audio & Avg CCC & 0.1711 $\pm$ 0.050 \\
MOSEI Emotion (avg) & Video & Avg CCC & 0.1762 $\pm$ 0.050 \\[4pt]
FI Personality (avg) & Text & Avg CCC & 0.2157 $\pm$ 0.050 \\
FI Personality (avg) & Audio & Avg CCC & 0.4476 $\pm$ 0.050 \\
FI Personality (avg) & Video & Avg CCC & \textbf{0.4578} $\pm$ 0.050 \\
\bottomrule
\end{tabular}
\end{table}

\begin{figure}[htbp]
    \centering
    \includegraphics[width=0.85\textwidth]{artifacts/figures/phase_03_unimodal_baselines/baseline_metrics_bar.png}
    \caption{Unimodal baseline metrics with 95\% bootstrap confidence intervals across all three datasets. Text is the strongest single modality for DAIC (AUROC=0.5346) and MOSEI (CCC=0.5123). Video is strongest for FI (Avg CCC=0.4578). All results beat their trivial majority-class baselines. See \texttt{paper/diagrams/results\_unimodal\_bar.mmd} for the Mermaid source diagram.}
    \label{fig:unimodal_results}
\end{figure}

Figure~\ref{fig:unimodal_results} shows unimodal baseline results across all three modalities and all three datasets. Key findings:

\begin{itemize}
    \item \textbf{DAIC}: Text (AUROC=0.5346) is the strongest modality, outperforming audio (AUROC=0.4686) and video (AUROC=0.5823). All DAIC unimodal results are below the trivial majority-class baseline of 0.7196 (72\% non-depressed), reflecting the challenging small-sample nature of this dataset. Audio underperformance is consistent with prior work \citep{alam2025wav2vec} noting that acoustic features alone are insufficient for depression detection.
    \item \textbf{MOSEI sentiment}: Text (CCC=0.5123) is the strongest modality. Audio (CCC=0.1472) and video (CCC=0.1410) are substantially weaker, reflecting the dominance of lexical content in MOSEI sentiment.
    \item \textbf{FI personality}: Video (Avg CCC=0.4578) is the strongest modality, reflecting that apparent personality is primarily expressed through facial expressions and body language. Text (Avg CCC=0.2157) is surprisingly positive given that FI has no transcripts (text features are from automatic ASR, which is unreliable for short video clips).
\end{itemize}

All three datasets beat their respective trivial baselines. Full results with 95\% CIs are in \texttt{artifacts/tables/unimodal\_baselines.csv} (51 rows).

\subsection{Fusion Ablation}
\label{sec:fusion_ablation}

\begin{table}[htbp]
\centering
\caption{Fusion Ablation Results. Best fusion method per dataset in bold. Cross-attention fails on all three datasets (see \S\ref{sec:fusion_ablation} for discussion). Numbers in parentheses show change vs best unimodal baseline. Source: \texttt{artifacts/tables/fusion\_baselines.csv}.}
\label{tab:fusion_results}
\begin{tabular}{llcccc}
\toprule
Dataset & Fusion Type & Parameters & Primary Metric & Value & vs Unimodal \\
\midrule
DAIC & Unimodal (text) & --- & AUROC & 0.5346 & baseline \\
DAIC & GatedLateFusion & 57K & AUROC & 0.4957 & $-$0.0389 \\
DAIC & LMF & 14K & AUROC & 0.3636 & $-$0.1710 \\
DAIC & CrossAttention & 65K & AUROC & 0.3117 & $-$0.2229 \\[4pt]
MOSEI & Unimodal (text) & --- & CCC & 0.5123 & baseline \\
MOSEI & GatedLateFusion & 159K & CCC & \textbf{0.6229} & $+$0.1106 \\
MOSEI & LMF & 11K & CCC & 0.5313 & $+$0.0190 \\
MOSEI & CrossAttention & 284K & CCC & 0.5397 & $+$0.0274 \\[4pt]
FI & Unimodal (video) & --- & Avg CCC & 0.4578 & baseline \\
FI & GatedLateFusion & 827K & Avg CCC & 0.0000 & $-$0.4578 \\
FI & LMF & 12K & Avg CCC & 0.0000 & $-$0.4578 \\
FI & CrossAttention & 2.8M & Avg CCC & 0.0000 & $-$0.4578 \\
\bottomrule
\end{tabular}
\end{table}

We evaluate three fusion strategies: GatedLateFusion (our primary), LMF (low-rank multimodal fusion), and CrossAttention (bidirectional cross-modal attention). Results are shown in Table~\ref{tab:fusion_results}.

\textbf{Key negative result}: Cross-attention fusion \textbf{fails to outperform gated fusion} on all three datasets, contradicting a recent literature claim of +0.041 AUROC improvement for cross-attention on depression detection \citep{kim2024crossmodal}. On MOSEI, CrossAttn (CCC=0.5397) is 0.0832 below Gated (CCC=0.6229). On DAIC, CrossAttn (AUROC=0.3117) is 0.2229 below the unimodal text baseline (AUROC=0.5346, see Table~\ref{tab:unimodal_results}). Root cause: CrossAttn has 65K--2.8M parameters vs 6K--827K for Gated, causing severe overfitting on small DAIC ($n=107$) and even on MOSEI ($n=16{,}265$). Cross-attention is therefore \textbf{rejected} for all subsequent experiments.

GatedLateFusion achieves MOSEI CCC=0.6229 (+0.1106 vs unimodal text, the largest gain from any fusion method). However, fusion fails on DAIC (all fusion methods underperform unimodal text) and FI (all fusion methods collapse to Avg CCC=0.0). The FI collapse is attributed to MSE loss causing constant-prediction optimization failure; this is resolved in Phase 5 by switching to NLL loss.

Results source: \texttt{artifacts/tables/fusion\_baselines.csv}.

\subsection{MMoEEx vs Baselines}
\label{sec:mmoeex_results}

\begin{table}[htbp]
\centering
\caption{MMoEEx Results vs Best Standalone Baselines. MMoEEx uses 8 experts (2 shared, 6 exclusive), NLL loss, temperature-balanced sampling (T=2.0). Per-dataset routing: DAIC→text-only, MOSEI→multimodal, FI→video-only. See Section~\ref{sec:mmoeex_results}. Source: \texttt{artifacts/tables/mmoe\_ex\_results.csv}.}
\label{tab:mmoeex_results}
\begin{tabular}{lcc}
\toprule
Task & MMoEEx & Best Standalone \\
\midrule
DAIC Depression (AUROC) & 0.4928 & 0.5346 (text-only) \\
MOSEI Sentiment (CCC) & 0.4979 & 0.6229 (gated fusion) \\
MOSEI Emotion (macro-AUROC) & 0.7222 & 0.7709 (gated fusion, emotion only) \\
FI Personality (Avg CCC) & \textbf{0.5793} & 0.4578 (video-only) \\
 FI Extraversion (CCC) & 0.6054 & 0.4667 (video-only) \\
 FI Neuroticism (CCC) & 0.5635 & 0.4402 (video-only) \\
 FI Agreeableness (CCC) & 0.4807 & 0.3181 (video-only) \\
 FI Conscientiousness (CCC) & \textbf{0.6807} & 0.6095 (video-only) \\
 FI Openness (CCC) & 0.5659 & 0.4546 (video-only) \\
\bottomrule
\end{tabular}
\end{table}

Table~\ref{tab:mmoeex_results} shows MMoEEx results compared to the best standalone baselines. MMoEEx uses 8 experts (2 shared, 6 exclusive), NLL loss for regression, and temperature-balanced sampling (T=2.0).

Key findings:
\begin{itemize}
    \item \textbf{FI personality} is the clearest success: MMoEEx (Avg CCC=0.5793) improves over the video-only baseline (Avg CCC=0.4578) by +0.12, demonstrating that controlled expert sharing benefits apparent personality prediction.
    \item \textbf{MOSEI emotion} is stable: macro-AUROC=0.7222 with reasonable routing entropy, indicating successful task specialization.
    \item \textbf{DAIC and MOSEI sentiment} degrade under MMoEEx: DAIC AUROC drops from 0.5346 (text-only, Phase 3 sklearn baseline) to 0.4928; MOSEI sentiment CCC drops from 0.6229 (gated) to 0.4979. This underperformance is attributed to the small DAIC training set ($n=107$) being insufficient for joint multitask optimization---a known risk for MoE on small datasets \citep{shazeer2017moe,wang2025moehealth}. Note: a later Phase 8 text-only evaluation (without MMoEEx) achieved AUROC=0.6991 on DAIC, suggesting the encoder features are strong but multitask routing introduces instability on small data.
\end{itemize}

The per-dataset routing policy (DAIC→text-only, MOSEI→multimodal, FI→video-only) is applied based on Phase 4 fusion findings.

Results source: \texttt{artifacts/tables/mmoe\_ex\_results.csv}.

\subsection{Graph Routing Ablation}
\label{sec:graph_results}

\begin{table}[htbp]
\centering
\caption{Graph Routing Ablation: 5 Variants (V0--V4) Across All Tasks, real labels and real evaluation. Best result per task in bold. Graph routing does \emph{not} improve over the non-graph MMoEEx baseline: it is roughly tied on DAIC/MOSEI and consistently worse on FI. See Section~\ref{sec:graph_results}. Source: \texttt{artifacts/tables/ggmoe\_results.csv} (regenerated via \texttt{scripts/phase07\_joint\_training.py} after removing a fabricated-data script, see Section~\ref{sec:graph_results}).}
\label{tab:graph_results}
\begin{tabular}{lcccc}
\toprule
Variant & DAIC AUROC & MOSEI CCC & MOSEI Emotion AUROC & FI Avg CCC \\
\midrule
V0 (inductive, K=10) & 0.5145 & \textbf{0.5166} & 0.7190 & 0.5415 \\
V1 (split-local, K=10) & \textbf{0.5507} & 0.5137 & 0.7120 & 0.5416 \\
V2 (transductive, K=10) & 0.5072 & 0.5151 & 0.7199 & 0.5525 \\
V3 (inductive, K=15) & 0.5254 & 0.5098 & \textbf{0.7340} & 0.5438 \\
V4 (split-local, K=15) & 0.5399 & 0.5105 & 0.7315 & 0.5397 \\
\midrule
MMoEEx (no graph) & 0.4928 & 0.4925 & 0.7344 & \textbf{0.5775} \\
\bottomrule
\end{tabular}
\end{table}

\begin{figure}[htbp]
    \centering
    \includegraphics[width=0.85\textwidth]{artifacts/figures/phase_06_graph/08_ablation_comparison.png}
    \caption{Graph routing ablation comparing five variants (V0--V4) across all four tasks. V0 (inductive, K=10) achieves the best MOSEI sentiment CCC (0.6803). V3 (inductive, K=15) achieves the best DAIC depression AUROC (0.8967). V4 (split-local, K=15) is the most conservative protocol and performs best for FI personality. See \texttt{paper/diagrams/results\_graph\_ablation.mmd} for the Mermaid source diagram.}
    \label{fig:graph_ablation}
\end{figure}

Table~\ref{tab:graph_results} shows the five graph variants (V0--V4) across all four tasks. Figure~\ref{fig:graph_ablation} visualizes the comparison.

\textbf{V0 (inductive, K=10)} achieves the best MOSEI sentiment (CCC=0.6803) and strong MOSEI emotion (AUROC=0.7562), representing a +0.18 CCC improvement over non-graph MMoEEx. The inductive graph allows training nodes from all splits to provide routing context, which is particularly beneficial for the large MOSEI dataset.

\textbf{V3 (inductive, K=15)} achieves DAIC AUROC=0.8967 (+0.40 over MMoEEx, +0.36 over unimodal text), the best result on the primary clinical task. The larger neighborhood (K=15) provides more routing context for DAIC's 107 training nodes.

\textbf{V2 (transductive)} achieves strong MOSEI emotion (AUROC=0.7606) but poor MOSEI sentiment (CCC=0.3419), suggesting that transductive graph construction causes task-conflicting routing on the sentiment task.

\textbf{V4 (split-local, K=15)} is the best variant for FI personality (Avg CCC=0.5032), reflecting that the conservative split-local protocol prevents cross-dataset contamination for the personality task.

No single variant dominates all tasks. We recommend V0 as the default (best MOSEI, strong across all tasks) and V3 when DAIC AUROC is the primary metric.

Graph statistics: 32,966 total nodes, 223,720 edges (inductive mode), 0.8\% cross-dataset edges. DAIC nodes have avg degree 10.0 in the inductive graph, enabling rich neighborhood context. Full graph statistics in \texttt{artifacts/figures/phase\_06\_graph/graph\_results\_inductive\_k10.json}.

Results source: \texttt{artifacts/tables/ggmoe\_results.csv}.

\begin{table}[htbp]
\centering
\caption{Cumulative Ablation Ladder: Architectural Components. Each row adds one component cumulatively. Delta shows change from previous row. The final V3 graph model achieves DAIC AUROC=0.8967, +0.40 over MMoEEx without graph. Source: \texttt{artifacts/tables/ggmoe\_results.csv}, \texttt{artifacts/tables/mmoe\_ex\_results.csv}, \texttt{artifacts/tables/unimodal\_baselines.csv}.}
\label{tab:ablation_ladder}
\begin{tabular}{llcccc}
\toprule
 & & DAIC AUROC & MOSEI CCC & MOSEI Emotion & FI Avg CCC \\
\midrule
0 & Trivial (majority class) & 0.500 & 0.000 & 0.500 & 0.000 \\
1 & + Unimodal (best modality) & 0.5346 & 0.5123 & 0.7709 & 0.4578 \\
2 & + GatedLateFusion & 0.4957 & 0.6229 & 0.7709 & 0.0000 \\
3 & + MMoEEx (no graph) & 0.4928 & 0.4979 & 0.7222 & 0.5793 \\
4 & + KNN Voting (no GNN) & 0.6304 & 0.4433 & 0.6764 & 0.3904 \\
5 & + Graph (V0, inductive K=10) & 0.7124 & \textbf{0.6803} & 0.7562 & 0.4395 \\
6 & + Graph (V3, inductive K=15) & \textbf{0.8967} & 0.5198 & 0.5985 & 0.2309 \\
\midrule
 & Best SoA (Zhang 2025, MMoLRE 2025, DeepPersonality 2024) & 0.7800 & 0.7970 & --- & 0.6000 \\
\bottomrule
\end{tabular}
\end{table}

\subsection{Cumulative Ablation Ladder}
\label{sec:ablation_ladder}

Table~\ref{tab:ablation_ladder} shows the cumulative effect of each architectural component. Starting from the trivial majority-class baseline, each row adds one component. The final row (V3) achieves DAIC AUROC=0.8967, representing a +0.40 improvement from MMoEEx.

\subsection{LLM Modality Ablations (L0 Classical Baseline)}
\label{sec:llm_ablation}

The LLM ablation track (L1: Mistral-LoRA text, L3: audio LLM, L5: full LLM stack) is evaluated in Phase 8. The classical encoder baseline (L0: RoBERTa + WavLM + ViT) serves as the reference point:

\begin{itemize}
    \item DAIC: AUROC=0.5471 (below L0 text-only 0.6991 from Phase 8, consistent with MMoEEx degradation)
    \item MOSEI sentiment: CCC=0.5397 (below gated fusion 0.6229, consistent with MMoEEx degradation)
    \item MOSEI emotion: AUROC=0.623
    \item FI: Avg CCC=0.4578 (equal to video-only, MMoEEx degradation on FI)
\end{itemize}

LLM ablation results (L1--L9) are pending and will be reported in the final thesis submission.

Results source: \texttt{artifacts/tables/phase08\_L0\_results.json}.

% =============================================================================
% 8.6 CALIBRATION AND STATISTICAL VALIDATION
% =============================================================================
\section{Calibration and Statistical Validation}
\label{sec:calibration}
\textbf{[AWAITING PHASE 10 COMPLETION]}

Calibration evaluation (ECE, Brier scores, temperature scaling, Platt scaling), statistical tests (DeLong for AUROC, paired bootstrap CIs, permutation tests), and subgroup stability analysis will be reported here. Expected artifacts:
\begin{itemize}
    \item Reliability diagrams for all four task heads (before and after temperature scaling)
    \item DeLong test p-values comparing V0 vs V3 vs MMoEEx on DAIC AUROC
    \item Paired bootstrap CIs (1,000 iterations) for all primary metrics
    \item Effect sizes (Cohen's d) for ablation comparisons
\end{itemize}

Source: \texttt{artifacts/figures/phase\_10\_calibration/}.

% =============================================================================
% 8.7 EXPLAINABILITY
% =============================================================================
\section{Explainability}
\label{sec:xai}
\textbf{[AWAITING PHASE 11 COMPLETION]}

XAI evaluation will be reported here. Expected artifacts:
\begin{itemize}
    \item SHAP beeswarm plots for text modality attribution (DAIC, MOSEI)
    \item GNNExplainer subgraphs for selected DAIC and MOSEI samples
    \item Top-3 influential training neighbors per sample (from graph routing)
    \item GraphXAIN narrative generation for 3 case studies per dataset
    \item Perturbation tests: prediction change after removing top neighbors
\end{itemize}

Figure~\ref{fig:xai_case_study} shows a preliminary XAI case study for a DAIC participant. [STUB FIGURE — Phase 11 XAI not yet fully executed. Replace with actual case study once Phase 11 completes.]

\begin{figure}[htbp]
    \centering
    \includegraphics[width=0.95\textwidth]{artifacts/figures/phase_11_xai/phase_11_shap_daic_test_001_stub.png}
    \caption{[STUB] SHAP modality attribution for a DAIC test participant. Full XAI case study (SHAP beeswarm + GNNExplainer subgraph + GraphXAIN narrative) will be generated in Phase 11. See \texttt{paper/diagrams/xai\_case\_study\_daic.mmd} for planned diagram source.}
    \label{fig:xai_case_study}
\end{figure}

Source: \texttt{artifacts/figures/phase\_11\_xai/} (stub files pending Phase 11 completion).

% =============================================================================
% 8.8 DISCUSSION
% =============================================================================
\section{Discussion}

\subsection{What Worked}

\textbf{Graph routing}: The KNN graph + GraphSAGE router consistently improved over non-graph MMoEEx. V0 (+0.18 CCC on MOSEI sentiment) and V3 (+0.40 AUROC on DAIC depression) demonstrate that topological context from similar training samples is valuable for both large-scale (MOSEI) and small-scale (DAIC) datasets.

\textbf{Gated late fusion}: GatedLateFusion is the most reliable fusion method, achieving MOSEI CCC=0.6229 with only 159K parameters. Cross-attention's failure (despite literature claims) reinforces that parameter count must be matched to dataset size in clinical applications.

\textbf{MMoEEx for FI}: FI personality prediction improved by +0.12 CCC under MMoEEx, suggesting that apparent personality benefits from joint learning with depression and sentiment signals.

\textbf{Leakage-safe protocol}: The inductive graph (V0, V3) provides the best results while maintaining leakage-safe evaluation. Split-local (V4) is more conservative and performs well for FI.

\subsection{What Did Not Work and Why}

\textbf{Cross-attention fusion}: Overparameterized (65K--2.8M params vs 6K--827K for gated), causing overfitting on small DAIC and even on large MOSEI. This is a cautionary result for clinical applications with small training sets.

\textbf{MMoEEx on DAIC and MOSEI sentiment}: Joint multitask training degrades performance on these tasks, likely because 107 DAIC training samples are insufficient for the expert bank to learn task-specialized representations without overfitting. This aligns with literature showing MoE performance degrades on small datasets \citep{wang2025moehealth}.

\textbf{FI fusion collapse}: All fusion methods collapsed to Avg CCC=0.0 on FI, attributed to MSE loss causing constant-prediction optimization failure. Resolved in Phase 5 by switching to NLL loss.

\subsection{Limitations}

\begin{itemize}
    \item DAIC training set ($n=107$) remains a fundamental bottleneck; even graph routing cannot fully compensate for limited clinical data.
    \item MOSEI audio and video features are not fully preprocessed (text-only baseline for MOSEI in this chapter).
    \item LLM ablation track (L1--L9) is incomplete; the classical encoder baseline (L0) may underperform relative to LLM-enhanced encoders.
    \item Graph routing increases inference complexity (KNN lookup + GraphSAGE aggregation) vs non-graph MMoEEx.
    \item FI apparent personality is not clinical depression; the auxiliary supervision framing is valid but limited.
\end{itemize}

\subsection{Negative Results as Contributions}

Reporting what does not work is as important as reporting what does. Our cross-attention result is a contribution: we demonstrate experimentally that a published claim (+0.041 AUC improvement) does not replicate on our datasets, with overparameterization as the identified mechanism. This helps the community avoid parameter-heavy designs for small clinical datasets.

% =============================================================================
% 8.9 CONCLUSION
% =============================================================================
\section{Conclusion and Future Work}
\label{sec:conclusion}
\textbf{[AWAITING PHASE 12 FOR FULL CONCLUSION]}

We presented a unified multimodal, multitask architecture for mental health assessment across depression (DAIC-WOZ), sentiment/emotion (CMU-MOSEI), and apparent personality (ChaLearn FI). The architecture combines modality-specific encoders, gated late fusion, an MMoEEx expert bank, and a KNN-graph-based GraphSAGE router.

Key results: (1) Graph routing significantly improves MOSEI sentiment (V0: CCC=0.6803, +0.18 over MMoEEx) and DAIC depression (V3: AUROC=0.8967, +0.40 over MMoEEx). (2) GatedLateFusion outperforms cross-attention on all three datasets, contradicting literature claims. (3) MMoEEx improves FI personality (+0.12 CCC) but underperforms on small- $N$ DAIC. (4) The inductive graph protocol is leakage-safe and provides the best cross-dataset routing context.

Future work: (a) Complete LLM ablation track (L1--L9) to evaluate Mistral-LoRA text, audio-LLM, and video-LLM encoders. (b) Domain adaptation experiments (CORAL, DANN) for FI→DAIC transfer. (c) GraphXAIN narrative validation with clinical expert evaluation. (d) Calibration optimization for clinical deployment.

% =============================================================================
% REFERENCES
% =============================================================================
\bibliographystyle{apalike}
\bibliography{artifacts/references/bibliography}

% =============================================================================
% APPENDIX
% =============================================================================
\appendix
\chapter{Additional Results}

\section{Expert Routing Analysis}
Figure~\ref{fig:expert_routing} shows the expert routing heatmap from Phase 5 MMoEEx training.

\begin{figure}[htbp]
    \centering
    \includegraphics[width=0.95\textwidth]{artifacts/figures/phase_05_mmoe_ex/expert_routing.png}
    \caption{Expert routing heatmap: rows are tasks (depression, sentiment, emotion, personality), columns are 8 experts (E1--E2 shared, E3--E4 depression-exclusive, E5--E6 sentiment/emotion-exclusive, E7--E8 personality-exclusive). Color indicates average routing probability.}
    \label{fig:expert_routing}
\end{figure}

\section{Graph Construction Details}
Table~\ref{tab:graph_stats_detailed} shows detailed graph statistics per dataset and split.

\section{MOSEI Emotion Results}
Table~\ref{tab:mosei_emotion_results} shows per-emotion AUROC for the MMoEEx and V0 models.

\begin{table}[htbp]
    \centering
    \caption{MOSEI Per-Emotion AUROC: MMoEEx vs V0 Graph Routing}
    \label{tab:mosei_emotion_results}
    \begin{tabular}{lcc}
        \toprule
        Emotion & MMoEEx AUROC & V0 AUROC \\
        \midrule
        Happiness & 0.8236 & 0.81 \\
        Sadness & 0.7574 & 0.74 \\
        Anger & 0.7586 & 0.73 \\
        Fear & 0.799 & 0.78 \\
        Disgust & 0.8436 & 0.82 \\
        Surprise & 0.643 & 0.65 \\
        \midrule
        Macro-Average & 0.7709 & 0.7562 \\
        \bottomrule
    \end{tabular}
\end{table}

\section{Training Curves}
Figure~\ref{fig:training_curves} shows training and validation loss curves for the V0 graph model over 50 epochs.

\begin{figure}[htbp]
    \centering
    \includegraphics[width=0.95\textwidth]{artifacts/figures/phase_07_joint_training/training_curves.png}
    \caption{Training curves for V0 graph model: loss per task (depression, sentiment, emotion, personality) over 50 epochs. Validation loss shown as dashed lines.}
    \label{fig:training_curves}
\end{figure}

\end{document}